% !TEX root = ../main.tex

\begin{acknowledgements}
  在撰写毕业论文期间, 我有幸得到中国科大苏州高研院 SCAI 实验室陈景润教授和孙羿斐博士的教导.
  两位老师不仅在研究思路和具体问题上给予我指导, 也在科研态度和人生规划上对我产生了深远影响.
  在此谨向他们的悉心教导和帮助致以诚挚的感谢.

  我对数学的热爱始于高中时期.
  在 17 岁以前, 我一直秉持实用主义的想法, 认为物理比数学更具有现实价值.
  那时我立志成为邓稼先那样推动我国科技进步的人.
  然而少年时期的心境常常会因一些偶然的契机而改变.
  17 岁生日时, 一位同学送给我一本数学科普读物\emph{万物皆数}, 由此引发了我对数学的兴趣.
  这本书聚焦于数学的实用价值和美学价值, 不仅打破了我对数学只是抽象工具且缺少实际应用的刻板印象, 也让我第一次真切地感受到数学的优雅.
  尽管它并没有直接促使我选择数学专业, 却在我心中埋下了接纳数学的种子.
  因此, 我也感谢这位同学和这本书对我的影响.

  我还要感谢父母和高中时期的老师们.
  由于学校课业繁重, 又有未来高考的压力, 我常常睡眠不足, 上课时有时会不由自主地打盹.
  我的班主任兼数学老师甄健华老师发现后, 并没有责备我, 而是与我和家长沟通, 一同寻找解决办法.
  最终, 父母在学校附近租下一套房子, 尽力保证我中午和晚上的休息.
  除班主任外, 其他科目的老师也都认真负责.
  语文老师奉行 ``学科素养教育'', 培养了我对文学和美学的热爱.
  英语老师授课严谨而认真, 在我考上中国科大的少创班后, 仍提醒我不能放下英语学习.
  作为物理课代表, 我与物理老师的互动较多, 她是一位认真负责而善良的老师.
  有时考虑到我们的作业较多, 她会减少作业量或延缓作业提交时间.
  在我离开高中时, 物理老师也因病需要治疗而暂别讲台, 所幸后来身体渐渐恢复, 又回到了教书育人的岗位.
  化学老师幽默风趣, 课堂上也常向我们分享成长经验, 在一定程度上塑造了我的人格.
  生物老师平易近人, 与学生关系融洽, 专业知识也十分丰富.

  中国科大的少年创新班给了我新的机遇.
  我的语文和英语成绩一直不算理想, 也曾常常担心自己会因偏科而无法进入喜欢的学校, 学习感兴趣的专业.
  中国科大的少创班给了我展示自己的机会.
  我十分感谢科大提供的平台, 使我能够心无旁骛地学习和科研.

  最后, 我要感谢我身边的朋友们.
  他们与我一起学习, 一起娱乐, 也在我遇到困难时伸出援手.
  本科生活中的许多乐趣, 都来自他们与我共同分享的快乐.
\end{acknowledgements}
